%% Monthly Update Slides -- March 2026
%% Metropolis theme (style/), green accent; content drawn from the paper
\documentclass[aspectratio=169,10pt]{beamer}

\usepackage{graphicx}
\usepackage{booktabs}
\usepackage[utf8]{inputenc}
\usepackage[T1]{fontenc}

\usetheme{metropolis}

%% Green accent for Metropolis
\definecolor{ornlgreen}{HTML}{2E7D32}
\setbeamercolor{normal text}{fg=black!90, bg=white}
\setbeamercolor{alerted text}{fg=ornlgreen}
\setbeamercolor{frametitle}{bg=ornlgreen, fg=white}
\setbeamercolor{title separator}{fg=ornlgreen}
\setbeamercolor{progress bar}{fg=ornlgreen}
\setbeamercolor{progress bar in head/foot}{fg=ornlgreen}
\setbeamercolor{progress bar in section page}{fg=ornlgreen}
\setbeamercolor{itemize item}{fg=ornlgreen}
\setbeamercolor{itemize subitem}{fg=ornlgreen!80!black}
\setbeamercolor{block title}{bg=ornlgreen!20, fg=ornlgreen!60!black}
\setbeamercolor{block body}{bg=ornlgreen!5}

%% Reduce body text to ~0.7 of nominal
\setbeamerfont{normal text}{size*={7pt}{8.4pt}}
\setbeamerfont{itemize/enumerate body}{size*={7pt}{8.4pt}}
\setbeamerfont{itemize/enumerate subbody}{size*={6.5pt}{8pt}}
\setbeamerfont{itemize/enumerate subsubbody}{size*={6pt}{7.5pt}}
\setbeamerfont{block title}{size*={8pt}{9.5pt},series=\bfseries}
\setbeamerfont{block body}{size*={7pt}{8.4pt}}
\setbeamerfont{frametitle}{size*={11pt}{13pt},series=\bfseries}
\setbeamerfont{title}{size*={14pt}{16pt},series=\bfseries}
\setbeamerfont{subtitle}{size*={10pt}{12pt}}
\setbeamerfont{author}{size*={9pt}{11pt}}
\setbeamerfont{institute}{size*={8pt}{10pt}}
\setbeamerfont{date}{size*={8pt}{10pt}}
\setbeamerfont{footnote}{size*={6pt}{7.2pt}}

\title{Explainable Machine Learning for Extreme\\Weather Event Prediction}
\subtitle{Long-Term Environmental Measurements at Oak Ridge}
\author{Jose Tupayachi \and Ziwei Liu \and Kevin Birdwell \and Xueping Li \and Xiao-Ying Yu}
\institute{Oak Ridge National Laboratory \and University of Tennessee, Knoxville}
\date{March 2026}

\begin{document}

%%--- Slide 1: Cover ---
\maketitle

%%--- Slide 2: Motivation & Contributions ---
\begin{frame}{Motivation \& Contributions}
\begin{columns}[T]
  \begin{column}{0.5\textwidth}
    \begin{block}{Problem}
      \begin{itemize}
        \item Extreme weather events (high winds, heavy precipitation,
              temperature excursions) threaten DOE infrastructure and
              operations at the Oak Ridge Reservation (ORR).
        \item Sub-hourly forecasts at $\le$\,1\,h horizon are needed
              to trigger protective actions.
        \item Standard ML models lack physical interpretability and
              are rarely benchmarked on long-term, in-situ tower data.
      \end{itemize}
    \end{block}
  \end{column}

  \begin{column}{0.5\textwidth}
    \begin{block}{Contributions}
      \begin{itemize}
        \item Curated \textbf{ORR-HiResMet}: $\sim$25 years of
              meteorological data from \textbf{6 towers}.
        \item End-to-end pipeline: extreme-event labelling,
              temporal/spatial/terrain features, ablation.
        \item Benchmark of \textbf{ML, DL, and CNN/TCN} models with
              Optuna-tuned hyper-parameters.
        \item \textbf{Explainability} via SHAP, permutation importance,
              and saliency linked to physical drivers.
        \item Cross-tower generalization \& computational efficiency.
      \end{itemize}
    \end{block}
  \end{column}
\end{columns}
\end{frame}

%%--- Slide 3: Data & Methods ---
\begin{frame}{Data \& Methodology}
\begin{columns}[T]
  \begin{column}{0.5\textwidth}
    \begin{block}{ORR-HiResMet Dataset}
      \begin{itemize}
        \item 6 meteorological towers across the ORR (TN, USA).
        \item Variables: wind speed/direction, T, RH, pressure,
              precipitation, solar radiation.
        \item Auxiliary: NSRDB irradiance, HRRR re-analysis,
              high-resolution terrain elevation.
        \item Quality control: duplicate detection, gap-filling,
              resampling to a uniform 1\,h grid.
      \end{itemize}
    \end{block}
    \begin{block}{Feature Engineering}
      \begin{itemize}
        \item \textbf{Temporal:} lags, rolling stats, gradients.
        \item \textbf{Spatial:} cross-tower deltas \& gradients.
        \item \textbf{Terrain:} elevation, slope, aspect.
      \end{itemize}
    \end{block}
  \end{column}

  \begin{column}{0.5\textwidth}
    \begin{block}{Models Benchmarked}
      \begin{itemize}
        \item \textbf{ML:} XGBoost, LightGBM, Random Forest.
        \item \textbf{DL:} LSTM, GRU.
        \item \textbf{CNN-based:} 1D-CNN, Temporal Convolutional Net.
        \item Hyper-parameters tuned with \textbf{Optuna}
              (TPE sampler, $>$100 trials per model).
      \end{itemize}
    \end{block}
    \begin{block}{Class Imbalance}
      \begin{itemize}
        \item Extreme events $<$1\% of samples.
        \item Focal loss + class-weighted sampling
              + threshold calibration.
      \end{itemize}
    \end{block}
  \end{column}
\end{columns}
\end{frame}

%%--- Slide 4: Results ---
\begin{frame}{Key Results}
\begin{columns}[T]
  \begin{column}{0.55\textwidth}
    \begin{block}{Event-Specific Performance}
      \begin{itemize}
        \item CNN/TCN best on high-wind events (highest F1, recall).
        \item LightGBM best on precipitation extremes (lowest FPR).
        \item GRU competitive on temperature excursions.
      \end{itemize}
    \end{block}
    \begin{block}{Ablation Study}
      \begin{itemize}
        \item Removing \textbf{spatial gradients} $\rightarrow$
              largest F1 drop ($-$8\,\%): cross-tower context matters.
        \item Removing \textbf{terrain features} $\rightarrow$
              consistent drop on wind events.
        \item Temporal-only baseline underperforms on rare events.
      \end{itemize}
    \end{block}
  \end{column}

  \begin{column}{0.42\textwidth}
    \begin{block}{Explainability \& Physics Linkage}
      \begin{itemize}
        \item SHAP highlights pressure tendency \& wind shear
              ahead of high-wind events.
        \item Saliency maps localize precipitation predictors
              to upwind towers (consistent with synoptic flow).
      \end{itemize}
    \end{block}
    \begin{block}{Cross-Tower \& Compute}
      \begin{itemize}
        \item Models trained on 5 towers transfer to the 6th
              with $\le$5\,\% F1 degradation.
        \item ML inference $<$1\,ms; DL/CNN $<$10\,ms per sample
              on a single GPU.
      \end{itemize}
    \end{block}
  \end{column}
\end{columns}
\end{frame}

%%--- Slide 5: Conclusion & Next Steps ---
\begin{frame}{Implications \& Next Steps}
\begin{columns}[T]
  \begin{column}{0.5\textwidth}
    \begin{block}{Implications}
      \begin{itemize}
        \item First \textbf{explainable, multi-tower} extreme-event
              benchmark on long-term ORR meteorological data.
        \item Spatial \& terrain features are not optional ---
              they are decisive for rare-event recall.
        \item Hybrid pipeline (ML for tabular extremes,
              CNN/TCN for sequence-driven events) gives
              the best operational trade-off.
      \end{itemize}
    \end{block}
  \end{column}

  \begin{column}{0.5\textwidth}
    \begin{block}{Conclusion}
      \begin{itemize}
        \item Open data \& code release supports reproducibility
              and DOE site-resilience studies.
        \item Framework extensible to other DOE sites via
              transfer learning.
      \end{itemize}
    \end{block}
    \begin{block}{Next Steps}
      \begin{itemize}
        \item Continual-learning updates as new tower data arrive.
        \item Fusion with HRRR/NSRDB for longer-horizon forecasts.
        \item Submitted to \emph{Journal of Big Data} --- March 2026.
      \end{itemize}
    \end{block}
  \end{column}
\end{columns}
\end{frame}

\end{document}
